\chapter{Uso e impacto de la IA y los LLM en el TFG}
\label{cap:capitulo11}

\vspace{0.25cm}

\section{Introducción}
\label{sec:segundaseccion}

En los últimos años, se han desarrollado numerosas herramientas de Inteligencia Artificial, concretamente aquellas relacionadas con los Modelos de Lenguaje de Gran Escala, también conocidos como Large Language Models (LLMs), lo que ha supuesto un cambio significativo en la manera de abordar una amplia diversidad de tareas, como por ejemplo, tareas de planificación, redacción y programación, en el caso de este TFG.

\vspace{0.25cm}

En el ámbito académico, estas herramientas se han convertido en un recurso habitual para los estudiantes, permitiendo mejorar la organización del trabajo o resolver dudas concretas. Sin embargo, el uso de estas herramientas se debe realizar de forma responsable, con el objetivo de garantizar que todo el contenido generado y proporcionado por las mismas sea comprendido, validado y adaptado dentro del trabajo o proyecto que se esté desarrollando.

\vspace{0.25cm}

En este capítulo se describen todos los usos que se han hecho de este tipo de herramientas, en qué medida, y en qué fases específicas han tenido mayor impacto dentro de este TFG.

\section{Herramientas utilizadas}
\label{sec:segundaseccion}

La única herramienta utilizada en este TFG basada en Inteligencia Artificial ha sido \href{https://chatgpt.com/es-ES/}{ChatGPT}, un modelo de lenguaje diseñado para ayudar a resolver y comprender tareas mediante la generación de texto.

\newpage

En el caso de este TFG, el uso de esta herramienta se ha limitado a la organización del trabajo, la estructuración de contenidos y la resolución de pequeñas dudas puntuales relacionadas con el desarrollo software. En ningún caso se ha utilizado como sustituto del trabajo técnico propio, sino como un recurso complementario que ha facilitado abordar fases del proyecto relacionadas con los casos anteriormente mencionados.

\section{Ámbitos de utilización}
\label{sec:segundaseccion}

Uno de los principales usos de \textit{ChatGPT} en este TFG ha sido como herramienta de apoyo para organizar y estructurar el código fuente desarrollado. Concretamente, para:

\begin{itemize}
\item Proponer formas de organizar el código de forma mucho más legible.
\item Resolver dudas puntuales relacionadas con la implementación de ciertas funciones.
\item Sugerir buenas prácticas de programación en Python para su posterior integración en ROS 2.
\end{itemize}

Este primer uso ha permitido obtener una estructuración del código clara y coherente, lo que ha facilitado su comprensión.

\vspace{0.25cm}

Otro de los usos de \textit{ChatGPT} en este TFG ha sido para definir la estructura de la memoria del mismo, además de servir como apoyo para:

\begin{itemize}
\item Definir un índice base para la memoria.
\item Organizar los capítulos de forma que los contenidos vistos en cada uno de ellos sean progresivos respecto a los anteriores.
\item Ajustar el tamaño de cada capítulo en función de la cantidad de trabajo llevada a cabo.
\end{itemize}

Este segundo uso ha permitido obtener una estructura alineada con el desarrollo de este TFG en lo que respecta a la memoria.

\vspace{0.25cm}

Y el último uso de \textit{ChatGPT} en este TFG ha sido relacionado con la planificación temporal de las distintas tareas llevadas a cabo, lo que ha permitido:

\newpage

\begin{itemize}
\item Estimar un tiempo límite aproximado para abordar cada tarea.
\item Organizar las tareas de forma secuencial y con una cantidad de trabajo incremental.
\item Ajustar la planificación en función del progreso desarrollado en todo momento.
\end{itemize}

Este tercer y último uso ha resultado útil para asegurar el cumplimiento de los plazos establecidos para todas las tareas desarrolladas en este TFG.

\section{Nivel de dependencia}
\label{sec:segundaseccion}

El uso de \textit{ChatGPT} en este TFG ha sido limitado y controlado en todo momento, de forma que todas sus funciones han sido únicamente auxiliares dentro del proceso de desarrollo.

\vspace{0.25cm}

Todo lo que respecta a decisiones técnicas, el diseño del sistema, la formulación del modelo cinemático, y la implementación y validación del mismo, han sido realizadas en su totalidad por el autor de este TFG. \textit{ChatGPT} se ha utilizado únicamente como una herramienta de apoyo para mejorar la organización del mismo y resolver dudas puntuales.

\vspace{0.25cm}

Por tanto, se puede afirmar que en ningún momento ha existido una dependencia total o directa de la herramienta, cuyo uso se ha limitado a mejorar la eficiencia del desarrollo de este TFG.

\section{Impacto en el desarrollo del TFG}
\label{sec:segundaseccion}

El impacto del uso de herramientas de Inteligencia Artificial en este TFG ha sido positivo, principalmente en aspectos relacionados con la organización y la planificación del mismo. Es por ello que hay que destacar tres grandes beneficios obtenidos:

\begin{itemize}
\item Mejora de la estructuración del trabajo tanto a nivel de código como a nivel de documentación.
\item Optimización del tiempo al facilitar la resolución de dudas concretas y la planificación de tareas.
\item Mayor claridad en la organización del contenido.
\end{itemize}

Sin embargo, este impacto se ha limitado única y exclusivamente a aspectos de apoyo, sin influir en ningún momento en el núcleo de este TFG.

\section{Limitaciones y consideraciones éticas}
\label{sec:segundaseccion}

El uso de cualquier herramienta basada en Inteligencia Artificial conlleva una serie de limitaciones que se deben tener en cuenta durante su uso.

\vspace{0.25cm}

En primer lugar, hay que tener en cuenta que este tipo de herramientas puede generar respuestas incorrectas o incompletas. Es por ello que toda la información proporcionada por \text{ChatGPT}  ha sido revisada y validada antes de ser utilizada o implementada en este TFG.

\vspace{0.25cm}

En segundo lugar, es importante mencionar que todos los usos de esta herramienta se han realizado de forma responsable, evitando en todo momento la generación de contenido sin haber obtenido previamente una comprensión del mismo, lo que ha resultado clave para garantizar la calidad y originalidad de este TFG.

\vspace{0.25cm}

Y por último, destacar que todos los usos de esta herramienta se han realizado respetando principios académicos y de autoría en todo momento, lo que ha permitido reflejar todo el esfuerzo dedicado y todos los conocimientos adquiridos a lo largo del grado.

\section{Conclusión}
\label{sec:segundaseccion}

El uso de \textit{ChatGPT} en este TFG se ha traducido en un apoyo únicamente complementario durante el desarrollo del mismo, centrándose en tareas de organización, estructuración y planificación, lo que ha permitido mejorar la calidad del proyecto y optimizar el tiempo de desarrollo de cada tarea, todo ello sin sustituir en ningún momento cualquier trabajo técnico propio.

\vspace{0.25cm}

Por ello, se puede concluir que el uso responsable de este tipo de herramientas puede aportar mucho más valor a todo el proceso de desarrollo en un proyecto académico, siempre que se mantenga un control riguroso sobre su uso y se garantice la comprensión, revisión y validación de toda la información proporcionada.
